\subsection{Solution to Problem 18: Curve on a Sphere}

\begin{solution}
The curve lies on the sphere $x^2 + y^2 + z^2 = 9$ of radius $3$. The given functions suggest using $z$ as the main parameter.

Let $t = z/3$, so $t \in [-1,1]$ as $z$ runs from $-3$ to $3$. The curve spirals three times around the $z$-axis, so the azimuthal angle should increase by $6\pi$ as $t$ goes from $-1$ to $1$. Take
\[
\theta = 3\pi t.
\]
At fixed height $z = 3t$, a point on the sphere has horizontal radius $\sqrt{9-z^2} = 3\sqrt{1-t^2}$.

Using the given functions $g(x)=\sqrt{9-x^2}$ and $f(x)=\cos(\pi x)$ with $x = z/3 = t$, a parameterisation is
\[
x(t) = 3\sqrt{1-t^2}\cos(3\pi t), \qquad
y(t) = 3\sqrt{1-t^2}\sin(3\pi t), \qquad
z(t) = 3t,
\qquad t \in [-1,1].
\]
Check that the curve stays on the sphere:
\begin{align*}
x^2 + y^2 + z^2
&= 9(1-t^2)\bigl(\cos^2(3\pi t)+\sin^2(3\pi t)\bigr) + 9t^2 \\
&= 9(1-t^2) + 9t^2 = 9.
\end{align*}
\end{solution}

\begin{takeaways}
\begin{itemize}
    \item Sphere constraint: $x^2+y^2+z^2=9$ with radius $3$
    \item Parameter $t=z/3$ makes the height vary linearly from $-3$ to $3$
    \item Spiralling: use $\theta=3\pi t$ and horizontal radius $3\sqrt{1-t^2}$
    \item The problem asks for "a possible set", not a unique solution
    \item The problem statement is not rigorous when it says "spirals 3 times" as this could be interpreted in different ways, relies entirely on the student's visual intuition
\end{itemize}
\end{takeaways}

\begin{remark}[A more rigorous formulation]
The original problem can be rephrased as follows.

A curve $\mathcal{C}$ lies on the sphere $x^2 + y^2 + z^2 = 9$. A particle moves along $\mathcal{C}$ such that its height is given by $z = t$ for $-3 \le t \le 3$.
\begin{enumerate}[label=\textbf{(\alph*)}]
    \item By considering only the $x$ and $y$ coordinates of the particle, show that its path traces a circle in the horizontal plane with a radius of $r(t) = \sqrt{9 - t^2}$.
    \item The particle rotates around the $z$-axis such that its angle $\theta$ increases linearly with $t$. If the particle completes exactly three full rotations between $t = -3$ and $t = 3$, find $\theta(t)$.
    \item Hence, write down the parametric equations for the curve $\mathcal{C}$.
\end{enumerate}

\medskip
\noindent\textbf{Outline of the solution.}
\begin{enumerate}[label=\textbf{(\alph*)}]
    \item With $z = t$ on the sphere, $x^2 + y^2 = 9 - t^2$, so at each height the particle lies on the circle $x^2 + y^2 = r(t)^2$ in the plane $z = t$, where $r(t) = \sqrt{9 - t^2}$.
    \item Three full rotations means $\theta(3) - \theta(-3) = 6\pi$. Since $\theta$ is linear in $t$, write $\theta(t) = at + b$; then $6a = 6\pi$, so $a = \pi$. One convenient choice is $\theta(t) = \pi t$ (any choice of $b$ shifts the starting angle but not the number of turns).
    \item In cylindrical coordinates, $x(t) = r(t)\cos\theta(t)$, $y(t) = r(t)\sin\theta(t)$, $z(t) = t$, giving
    \[
    x(t) = \sqrt{9 - t^2}\cos(\pi t), \qquad
    y(t) = \sqrt{9 - t^2}\sin(\pi t), \qquad
    z(t) = t,
    \qquad -3 \le t \le 3.
    \]
    This agrees with the parameterisation above after the substitution $t \mapsto z/3$.
\end{enumerate}
\end{remark}
